% Halaman Sampul Tesis (Lampiran 6 - Panduan UDINUS)
\newgeometry{left=4cm,right=2.5cm,top=2.5cm,bottom=2.5cm}
\thispagestyle{empty}
\begin{center}
\addcontentsline{toc}{chapter}{HALAMAN SAMPUL}
% Jarak vertikal dari margin atas ke blok judul.
% Ubah 1.5cm untuk menggeser seluruh konten ke atas (nilai lebih kecil) atau ke bawah (nilai lebih besar).
\vspace*{0.5cm}
{\large\bfseries\MakeUppercase{\ThesisType}}\\[1.0cm]
% Ukuran font judul dan program studi mengikuti default dokumen (12pt).
% Tambahkan perintah ukuran seperti \Large atau \large di sini jika ingin mengubahnya.
{\large\bfseries\MakeUppercase{\TitleID}}\\[1.0cm]

Oleh:\\[0.4cm]
{\bfseries\MakeUppercase{\AuthorName}}\\[0.1cm]
{\bfseries \StudentID}\\[1.0cm]

{\bfseries Tesis diajukan sebagai salah satu syarat\\
untuk memperoleh gelar\\
\DegreeName}\\[1.0cm]

% Logo universitas (diameter ~4cm)
\IfFileExists{\UniversityLogo}{\includegraphics[width=4cm]{\UniversityLogo}\\[1.0cm]}{\fbox{\parbox[c][4cm][c]{4cm}{\centering LOGO}}\\[1.0cm]}

\vfill
% Nama program studi, fakultas, dan universitas (ukuran default 12pt).
{\large\bfseries\MakeUppercase{\ProgramName}}\\
{\large\bfseries\MakeUppercase{\FacultyName}}\\
{\large\bfseries\MakeUppercase{\UniversityName}}\\
{\large\bfseries\MakeUppercase{\City}}\\
{\large\bfseries\Year}
\end{center}
\clearpage
\restoregeometry
